\documentclass[12pt]{article}
\usepackage[utf8]{inputenc}
\usepackage[T1]{fontenc}
\usepackage{hyperref}
\title{Samburu Language Resources for K-12 Curriculum Development}
\author{Compiled by Claude Sonnet 5.5 (claude-sonnet-5-5) for Andr\'as Kornai}
\date{}
\begin{document}
\maketitle

\section{Introduction}
Samburu (ISO 639-3 \texttt{saq}; Glottocode samb1315) is a Nilotic (Maa) language of northern Kenya with roughly 240,000 speakers. Written material is limited. The Bible and a few dictionaries and orthography documents are the main printed or online resources found.

\section{Primary Educational Resources}
I found no Samburu primers or textbooks. Kenya's language-in-education policy prescribes the learners' mother tongue in Grades 1--3 in rural areas, with English from Grade 4. I found no Samburu-specific implementation data.

\section{Linguistic and Reference Materials}
\begin{itemize}
\item Karani, Kotikash and Sentero, \emph{A Unified Standard Orthography for Maa Languages: Kenya and Tanzania} (CASAS, 2014), covering Maasai, Samburu and related varieties: \url{https://repository.udsm.ac.tz/items/81c696c8-a708-4db7-b974-406794305051/full}
\item Maa Language Project, University of Oregon (Doris Payne), with a preliminary online dictionary covering Samburu and Kenyan Southern Maasai: \url{https://pages.uoregon.edu/maasai/}
\item Carol Fabert, \emph{The Samburu-English Travel Dictionary} (Wamba, Kenya, 2022): \url{https://sambururesearch.com/the-english-samburu-travel-dictionary/}
\item Dictionary methodology notes: \url{https://sambururesearch.com/2022/02/16/dictionary-methodology/}
\end{itemize}

\section{Religious and Community Resources}
\begin{itemize}
\item News report on the launch of the Samburu New Testament (\emph{Lning'o Ng'ejuk}, Maralal, December 2019), accompanied by a dictionary and an audio version; the Old Testament was still in progress: \url{https://www.the-star.co.ke/news/2019-12-15-joy-as-bible-translated-in-samburu-launched}
\item Kenya News Agency report: \url{https://www.kenyanews.go.ke/samburu-language-bible-launched/}
\item Digital Bible Society listing: \url{https://dbs.org/bibles/SAQWBT}
\end{itemize}

\section{Digital and Audio Resources}
An audio version of the New Testament was reported at the 2019 launch; I did not see a direct link to it.

\section{Curriculum Development Context}
Samburu is stable and widely spoken, but there are almost no published texts apart from the Bible and its associated dictionary, and earlier spelling conventions were not harmonised. The unified Maa orthography and the Bible translation are the natural starting points for a standard written form. Teaching texts would have to be newly written, and no Samburu-specific policy or school data was found.

\end{document}